\chapter{\code{IngestAdapter} Trait}
\label{ch:ingest-adapter}

\begin{note}
Stub. The P2-E refactor introduced a higher-level wrapper around
\code{LogSource} + \code{LogParser}: new source authors implement
\code{IngestAdapter} in a child crate rather than editing core
files.
\end{note}

\section{Contract}

\begin{lstlisting}[style=rust]
#[async_trait]
pub trait IngestAdapter: Send + Sync {
    fn adapter_id(&self) -> &str;
    fn tables(&self) -> &[&'static str];
    async fn stream_triples(
        &self,
        table: &str,
        since: Watermark,
        cancel: CancellationToken,
    ) -> Result<TripleBatchStream, SourceError>;
}

pub struct TripleBatch {
    pub triples:    Vec<ParsedTriple>,
    pub watermark:  Watermark,
    pub fetched_at: i64,
}
\end{lstlisting}

\section{Generic implementation}

\code{LogSourceAdapter} wraps any \code{LogSource} + \code{LogParser}
into an \code{IngestAdapter}. The 4 existing sources (Log
Analytics, Defender XDR, Data Lake, AAD export) can be exposed as
adapters with zero source-code change; the orchestrator still holds
them in concrete form today for backward compatibility, but the
seam is open.

\section{MVCC propagation}

\code{TripleBatch.fetched\_at} carries the upstream fetch timestamp,
which is exactly what the writer needs to call
\code{CompactRelation::with\_ingested\_at} on each inserted edge ---
the third leg of the P3 MVCC story (Chapter~\ref{ch:mvcc}).
Wiring that plumbing end-to-end is a deferred follow-up.

\begin{inthecode}
\filepath{graph\_hunter\_core/src/ingest\_adapter.rs} --- trait +
\code{LogSourceAdapter}.
\end{inthecode}
